\section{Model Parameters and Experimental Details}
\label{app:model_parameters_experimental_details}
\providecommand{\modelappendixfigdir}{runs/appendix_figures/mean_pooling_model}

This appendix reports the configuration used for the no-overlap LEAVES pretraining
of the mean-pooling model.  The LEAVES manifest excludes duplicated StarEmbed
targets and contains 1,056,785 light curves.
The active pretraining views are raw light curve, periodogram, and phase-folded
light curve.  The group-fusion branch and projection head are disabled.

\subsection{LEAVES--StarEmbed Deduplication}
\label{app:leaves_starembed_deduplication}

The LEAVES pretraining set and the StarEmbed benchmark are both derived from
ZTF light curves, so target-level deduplication is required before using
StarEmbed as a downstream evaluation benchmark.  The filtering is performed at
the astronomical-source level rather than by removing all ZTF targets.  In
particular, only LEAVES light curves whose target is crossmatched to a
StarEmbed source are removed; all other ZTF light curves remain in the LEAVES
pretraining corpus.

We first construct the set of StarEmbed evaluation targets from the preprocessed
StarEmbed splits used in this work: train, validation, test, and anomaly.  Each
StarEmbed preprocessed file name contains the Pan-STARRS source identifier, and
these identifiers are collected from the split manifests.  The resulting
Pan-STARRS source set contains 40,716 unique StarEmbed targets.  We then query
the StarEmbed crossmatch table, which provides Pan-STARRS source identifier,
Gaia source identifier, right ascension, and declination.  Restricting this
table to the collected StarEmbed sources gives 41,045 catalog rows, 40,589
unique Gaia identifiers, and 40,590 unique coordinate entries with valid
coordinates.

LEAVES targets are matched to this StarEmbed source catalog in two passes.  The
first pass checks whether a LEAVES target name explicitly contains a Gaia
identifier of the form \texttt{GaiaID<source\_id>}.  If that Gaia identifier is
present in the StarEmbed crossmatch table, the LEAVES target is removed.  The
second pass is a positional crossmatch for LEAVES target names that encode sky
coordinates.  We parse the target coordinate string as
\texttt{JHHMMSS.s+DDMMSS.s} or \texttt{JHHMMSS.s-DDMMSS.s}, convert the right
ascension and declination to degrees, and represent both LEAVES and StarEmbed
positions as three-dimensional unit vectors on the celestial sphere.  A
\texttt{cKDTree} is built over the StarEmbed unit vectors, and a LEAVES target
is removed if it has any StarEmbed neighbor within 2 arcsec.  The chord radius
used by the tree query is
\[
    r = 2\sin\left(\frac{\theta}{2}\right),
    \qquad \theta = 2~\mathrm{arcsec},
\]
with $\theta$ converted to radians.  This spherical-vector query avoids
right-ascension wrap-around issues and applies the same angular threshold over
the full sky.

The final no-overlap manifest is written by scanning the original LEAVES
manifest and retaining only targets that fail both the Gaia-identifier match and
the 2-arcsec positional match.  In our run, the original LEAVES manifest
contains 1,099,827 light curves.  The coordinate parser is applicable to
965,231 LEAVES target names; 43,042 light curves are removed by the positional
crossmatch, and no additional light curves are removed by the Gaia-identifier
pass.  The resulting no-overlap LEAVES manifest contains 1,056,785 light
curves.  This manifest is used for all final LEAVES pretraining and for the
re-pretraining of comparison models.

\begin{table*}[t]
    \centering
    \caption{Backbone parameterization for the no-overlap full model.}
    \label{tab:appendix_model_parameters}
    \begin{tabular}{l l l}
        \specialrule{2pt}{1pt}{1pt}
        Component & Setting & Value \\
        \specialrule{2pt}{1pt}{1pt}
        Raw-view encoder & Transformer depth & 4 layers \\
        Raw-view encoder & Hidden width & 256 \\
        Raw-view encoder & Attention heads & 4 \\
        Raw-view encoder & Position encoding & C-RoPE over observation time \\
        Raw-view encoder & Numeric embedding & Error-aware Gaussian embedding, 2,048 bins \\
        Raw-view encoder & Pooling & Masked mean pooling \\
        Raw-view encoder & Output dimension & 128 \\
        Raw-view encoder & Parameters & 3,713,664 \\
        \hline
        Periodogram-view encoder & Transformer depth & 4 layers \\
        Periodogram-view encoder & Hidden width & 256 \\
        Periodogram-view encoder & Attention heads & 4 \\
        Periodogram-view encoder & Position encoding & C-RoPE over $\log_{10} P$ \\
        Periodogram-view encoder & Numeric embedding & Log-power embedding, 2,048 bins \\
        Periodogram-view encoder & Pooling & Masked mean pooling \\
        Periodogram-view encoder & Output dimension & 128 \\
        Periodogram-view encoder & Parameters & 3,713,664 \\
        \hline
        Phase-folded-view encoder & Transformer depth & 4 layers \\
        Phase-folded-view encoder & Hidden width & 256 \\
        Phase-folded-view encoder & Attention heads & 4 \\
        Phase-folded-view encoder & Position encoding & C-RoPE over normalized phase \\
        Phase-folded-view encoder & Numeric embedding & Error-aware Gaussian embedding \\
        Phase-folded-view encoder & Pooling & Masked mean pooling \\
        Phase-folded-view encoder & Output dimension & 128 \\
        Phase-folded-view encoder & Parameters & 3,713,664 \\
        \hline
        Backbone & Enabled views & Raw, periodogram, phase-folded \\
        Backbone & Disabled components & Group fusion and projection head \\
        Backbone & Concatenated representation & 384 dimensions \\
        Backbone & Total parameters & 11,140,992 \\
        \specialrule{2pt}{1pt}{1pt}
    \end{tabular}
\end{table*}

\begin{table*}[t]
    \centering
    \caption{Pretraining and detached-probe optimization settings for the no-overlap full model.}
    \label{tab:appendix_training_details}
    \begin{tabular}{l l}
        \specialrule{2pt}{1pt}{1pt}
        Setting & Value \\
        \specialrule{2pt}{1pt}{1pt}
        Pretraining objective & LeJEPA only, weight 1.0, $\lambda=0.02$; CLIP loss disabled \\
        Regularization & SIGReg with 17 knots, maximum $t=3.0$, projection dimension 256 \\
        Views in loss & Raw, periodogram, phase-folded \\
        Numeric value bins & 2,048 bins for raw magnitudes and periodogram log-power values \\
        Raw magnitude range & $[-4, 4]$ with error-aware Gaussian numeric embedding \\
        Periodogram log-power range & $[-6, 0]$ \\
        Batch size & 128 per process; 8 DDP processes; effective global batch size 1,024 \\
        Training length & 10,000 optimizer steps; gradient accumulation 1 \\
        Backbone optimizer & AdamW, learning rate $2\times10^{-4}$, weight decay 0.01 \\
        Backbone AdamW parameters & $\beta_1=0.9$, $\beta_2=0.999$, $\epsilon=10^{-8}$ \\
        Backbone gradient clipping & Global norm 1.0 \\
        Mixed precision & Enabled \\
        Data split & Test ratio 0.1, seed 1234 \\
        Logging and checkpoints & Log every 10 steps; save every 10,000 steps \\
        Detached-probe features & Concatenated detached embeddings from the three active views, dimension 384 \\
        Detached-probe optimizer & AdamW, learning rate $10^{-3}$, weight decay 0 \\
        Detached-probe losses & Cross-entropy for 7-class and 10-class labels; labels $<0$ ignored \\
        \specialrule{2pt}{1pt}{1pt}
    \end{tabular}
\end{table*}

\begin{figure*}[t]
    \centering
    \includegraphics[width=0.92\textwidth]{\modelappendixfigdir/detached_probe_training_curves.pdf}
    \caption{Training curves for the detached embedding probes during no-overlap pretraining.
    The pale traces show batch-level logged values every 10 steps, and the solid traces show
    a moving average.  The probes are optimized on detached concatenated embeddings and do not
    backpropagate into the backbone.  At the final logged step, the 7-class probe reaches
    90.92\% batch accuracy with loss 0.270, and the 10-class probe reaches 89.45\% batch
    accuracy with loss 0.317.}
    \label{fig:appendix_detached_probe_training_curves}
\end{figure*}

\subsection{Model Training Details}
\label{app:model_training_details}

All LEAVES pretraining experiments in the final comparison use the no-overlap
LEAVES manifest, which removes duplicated StarEmbed targets before pretraining.
This prevents the frozen StarEmbed evaluation split from appearing in the
self-supervised pretraining corpus.  Unless otherwise stated, models are trained
for 10,000 optimizer steps with seed 1234, batch size 128 per process, eight DDP
processes, gradient accumulation 1, gradient clipping at global norm 1.0, and
AdamW with learning rate $2\times10^{-4}$, weight decay 0.01, betas
$(0.9,0.999)$, and $\epsilon=10^{-8}$.  The detached probes use concatenated
frozen embeddings, learning rate $10^{-3}$, and no weight decay.

\begin{itemize}
    \item \textbf{Our model}: Our main model uses three numeric views: the raw
    light curve, the Lomb--Scargle periodogram, and the phase-folded light
    curve.  The group branch and projection head are disabled.  Each view
    encoder has width 256, four transformer blocks, four attention heads, and a
    128-dimensional output embedding.  The raw and phase-folded views use
    time/phase C-RoPE and error-aware numeric embeddings; the periodogram view
    uses C-RoPE over periodogram coordinates.  The final sequence
    representation is obtained by masked mean pooling over valid transformer
    tokens.  The pretraining objective is LeJEPA with weight 1.0 and
    $\lambda=0.02$, together with SIGReg using 17 knots, maximum $t=3.0$, and
    projection dimension 256.  CLIP loss is disabled for the main model.  The
    three 128-dimensional view embeddings are concatenated for detached
    probing and StarEmbed extraction, giving a 384-dimensional representation.
    The backbone has 11,140,992 trainable parameters.

    \item \textbf{Astromer-2}: Astromer-2 is re-pretrained on the same
    no-overlap LEAVES manifest rather than using the original public checkpoint.
    To make the comparison size-matched, its encoder is resized to match the
    encoder scale used by our model: $d_{\mathrm{model}}=256$, four
    transformer layers, four attention heads, feed-forward dimension 1024,
    dropout 0.1, and RoPE base 10000.  The Astromer-2 masking configuration
    uses probed fraction 0.5, masked fraction 0.3, and random fraction 0.1.
    The input sequence length is capped at 1000 for LEAVES pretraining, with
    period search range $[0.006944444,2000]$ over 1000 grid points.  The model
    uses the shared 10,000-step AdamW training recipe above, with batch size
    128, train fraction 0.9, point-drop augmentation probability 0.1, and
    maximum dropped-point fraction 0.10.  StarEmbed extraction uses sequence
    length 200, z-score normalization, batch size 256, and the final
    checkpoint.

    \item \textbf{FALCO}: FALCO is also re-pretrained from scratch on the same
    no-overlap LEAVES manifest and is resized to the encoder scale used by our
    model:
    $d_{\mathrm{model}}=256$, four transformer layers, four attention heads,
    feed-forward dimension 1024, dropout 0.1, and RoPE base 10000.  Error bars
    are included in the input representation, the Huber loss uses
    $\delta=1.0$, the maximum accepted error bar is 5.0, and times are
    normalized to start at zero.  It uses the same 10,000-step AdamW recipe,
    batch size 128, train fraction 0.9, mixed precision, gradient clipping at
    1.0, and point-drop augmentation probability 0.1 with maximum dropped-point
    fraction 0.10.  StarEmbed features are extracted from the final
    checkpoint with sequence length 200, batch size 256, and mean pooling.

    \item \textbf{Ablation models}: The ablation suite is trained with the same
    no-overlap data, masked mean pooling, encoder size, optimizer, and training
    length as the main model unless the ablated component changes the active
    views or objective.  We train one-component ablations that remove C-RoPE,
    remove the error-aware numeric embedding, remove the raw-view branch, remove
    the phase-folded-view branch, remove the GLS periodogram branch, or replace
    the LeJEPA objective with the symmetric CLIP objective.  These ablations
    isolate whether the gains come from temporal rotary encoding, uncertainty
    modeling, individual time-series views, or the multi-view predictive loss.
\end{itemize}

\section{Domain-specific adaption for PYRREGULAR dataset}
\label{app:pyrregularadaption}

To further test whether the proposed framework can be adapted beyond
astronomical light curves, we evaluate a general irregular-time-series variant
on PYRREGULAR.  PYRREGULAR differs from StarEmbed and LEAVES in both data
semantics and observation structure: each sample can contain multiple sensor or
clinical channels, the time axis is not necessarily astronomical observing
time, and the target labels describe heterogeneous domains such as clinical
state, human activity, animal movement, or trajectory type.  We therefore use a
domain-adapted configuration instead of assuming that the astronomy-optimized
model is universally optimal.

The domain-adapted PYRREGULAR run is initialized from the three-view LEAVES
model trained with attention-statistics pooling and a combined LeJEPA--CLIP
objective.  The three active self-supervised views are the raw sequence, the
periodogram, and the phase-folded sequence.  The raw and phase-folded views use
error-aware numeric embeddings and C-RoPE over time or phase, and the
periodogram view uses C-RoPE over periodogram coordinates.  Unlike the
astronomy-facing no-overlap model used for the main StarEmbed benchmark, this
variant uses attention-statistics pooling and keeps the symmetric CLIP term
between raw--periodogram and raw--phase-folded embeddings.  In the stored run
used here, the group-view fusion branch is disabled; the architecture supports
such metadata-level fusion, but it is not part of the reported PYRREGULAR
numbers.

\begin{table*}[t]
    \centering
    \caption{Experimental differences between the PYRREGULAR-adapted variant
    and the no-overlap astronomy model used for the main StarEmbed experiments.}
    \label{tab:pyrregular_adaptation_setup}
    \begin{tabular}{l l l}
        \specialrule{2pt}{1pt}{1pt}
        Setting & PYRREGULAR-adapted variant & Astronomy no-overlap model \\
        \specialrule{2pt}{1pt}{1pt}
        Source pretraining data & Original LEAVES manifest & No-overlap LEAVES manifest \\
        Active views & Raw, periodogram, phase-folded & Raw, periodogram, phase-folded \\
        Encoder scale per view & 4 layers, width 256, 4 heads & 4 layers, width 256, 4 heads \\
        View output dimension & 128 & 128 \\
        Pooling & Attention-statistics pooling & Masked mean pooling \\
        Group-view fusion & Disabled in reported run & Disabled \\
        Pretraining objective & LeJEPA + symmetric CLIP + SIGReg & LeJEPA + SIGReg \\
        CLIP pairs & Raw--periodogram, raw--phase-folded & Not used \\
        Projection head & Disabled & Disabled \\
        Trainable parameters & 12,126,339 & 11,140,992 \\
        PYRREGULAR SFT length & 30 epochs & 10 epochs \\
        PYRREGULAR SFT hardware & 1 GPU & 8 GPUs \\
        SFT mode & Full-model fine-tuning & Full-model fine-tuning \\
        SFT value scaling & None & None \\
        SFT time strategy & Relative time & Relative time \\
        Maximum points per series & 1000 & 1000 \\
        Downstream heads & Logistic, MLP, $k$-NN & Logistic, $k$-NN \\
        \specialrule{2pt}{1pt}{1pt}
    \end{tabular}
\end{table*}

\begin{table*}[t]
    \caption{PYRREGULAR macro-F1 comparison between the previous SOTA, the
    astronomy-facing no-overlap model, and the PYRREGULAR-adapted variant.
    The detailed meaning of dataset and previous SOTA method names follows the
    original PYRREGULAR paper \cite{spinnato2026pyrregularunifiedframeworkirregular}.
    For the PYRREGULAR-adapted variant, only Linear and $k$-NN heads are
    considered here.}
    \label{tab:PYRREGULAR_domain_adaptation}
    \centering
    \begin{tabular}{l|lc|lc|lc}
        \hline
        Dataset
        & Prev. SOTA & Prev. F1
        & No-overlap head & No-overlap F1
        & Adapted head & Adapted F1 \\
        \hline
        ABF & NCDE   & 0.41 & Linear & 0.37 & Linear & 0.55 \\
        AN  & ROCKET & 0.90 & $k$-NN & 0.58 & $k$-NN & 0.52 \\
        PGE & BRITS  & 0.78 & Linear & 0.78 & $k$-NN & 0.78 \\
        GS  & NCDE   & 0.52 & Linear & 0.20 & Linear & 0.20 \\
        LPA & BORF   & 0.73 & Linear & 0.82 & Linear & 0.77 \\
        MI3 & RIFC   & 0.56 & Linear & 0.69 & Linear & 0.61 \\
        PAM & ROCKET & 0.66 & Linear & 0.49 & Linear & 0.74 \\
        P12 & RIFC   & 0.63 & Linear & 0.58 & Linear & 0.61 \\
        P19 & LGBM   & 0.75 & $k$-NN & 0.50 & $k$-NN & 0.57 \\
        SE  & RIFC   & 0.82 & Linear & 0.67 & Linear & 0.67 \\
        TA  & LGBM   & 0.77 & Linear & 0.35 & Linear & 0.36 \\
        VE  & BORF   & 0.87 & Linear & 0.84 & $k$-NN & 0.64 \\
        \hline
    \end{tabular}
\end{table*}

For PYRREGULAR adaptation, each multivariate sample is decomposed into
univariate channels.  During dataset-specific self-supervised fine-tuning, the
non-test split channels are converted into the same raw, periodogram, and
phase-folded views used by the source model.  The model is then fine-tuned with
the original self-supervised objective of the source configuration, using
dataset-specific value ranges, time spans, and periodogram bounds estimated
from the training split.  During evaluation, each channel is encoded by the
raw-view encoder, channel embeddings are averaged into one sample embedding,
and lightweight classifiers are trained on the PYRREGULAR non-test split and
evaluated on the official test split.  This protocol separates the
self-supervised domain adaptation step from the supervised classifier used to
measure representation quality.

The comparison highlights that the domain-adapted configuration and the
astronomy-facing configuration differ in more than only training length.  The
PYRREGULAR-adapted variant keeps the contrastive alignment term and uses a
richer attention-statistics pooling operator, whereas the astronomy-facing
model removes CLIP and uses masked mean pooling because that combination gives
stronger StarEmbed performance after removing duplicated StarEmbed targets from
LEAVES.  The PYRREGULAR run also evaluates an MLP head in addition to logistic
regression and $k$-NN, while the 10-epoch no-overlap run evaluates only
logistic regression and $k$-NN.

With dataset-specific adaptation, the attention-statistics/LeJEPA--CLIP variant
obtains a mean best-head macro F1 of 58.74 across the 12 PYRREGULAR datasets,
compared with 57.26 for the 10-epoch no-overlap mean-pooling model.  The gains
are concentrated on several heterogeneous non-astronomical tasks: for example,
the adapted variant reaches macro F1 73.85 on PAMAP2, 60.78 on PhysioNet 2012,
58.82 on PhysioNet 2019, and 55.47 on ABF, while the corresponding no-overlap
mean-pooling values are 49.17, 58.00, 49.80, and 36.87.  Conversely, the
astronomy-facing model remains stronger on some individual datasets, including
MIMIC-III, LDFPA, Animals, and Vehicles.  This contrast supports our central
claim that time-series representation learning is domain-dependent: the
self-supervised multi-view framework is reusable, but the optimal choice of
pooling mechanism, objective terms, and adaptation protocol depends on the
semantics of the target domain.
